\section{Forecast-readout protocol and outcomes}
\label{app:forecast}
The three gate models stay frozen. Up to 16,384 training contexts and 8,192 evaluation contexts are
sampled without replacement with seeds 20260919 and 20260920, using matching context IDs across
models. Past context may not cross a reset. Forecasts use the original recorded outgoing actions,
with horizons 1, 2, 4, and 8. Parameters and BatchNorm statistics must remain unchanged during
extraction, and cached encodings are checked against the public rollout implementation.

The three feature domains are real future encoding, autoregressive forecast, and copied current
encoding. Each domain has 512 coordinates. An affine ridge read-out uses mean squared error plus
$10^{-3}\|W\|^2$, with unpenalized intercept and no hyperparameter search. Both declared scaler
analyses are retained: current-training-latent mean/standard deviation shared across domains; and
source-domain training mean/standard deviation for each horizon. Both use a fixed $10^{-6}$ scale
floor, and every transfer applies its source-fitted scaler. No evaluation data determine a scaler,
regularizer, checkpoint, or reporting condition.

Cube coordinates are divided by 63 for fitting and reported in pixels. Wall position is fitted and
evaluated only on truth-present rows; the absent-wall sentinel is never regressed as a coordinate.
Presence reports class prevalence, AUROC, average precision, and threshold accuracy/balanced
accuracy, with undefined one-class metrics omitted. Every task includes a constant training-label-mean
baseline and sample counts. A task with fewer than 64 fitting rows is marked insufficient.

The archive stores reset frames after terminal transitions, so physical-state read-outs use a
horizon-specific continuing-episode mask. In particular, one-step evaluation does not select for
eight-step survival. The hidden-gap negative control instead retains every pre-reveal context,
including later resets, because future-survival selection could reveal the gap through the recorded
actions. The gap label refers to the original current wall. Chance forecasting of this unavailable
bit is compatible with correct treatment of uncertainty.

All fitted weights, scalers, selected IDs, labels, masks, per-example predictions, and provenance
hashes are saved. Overlapping windows are not independent samples; any uncertainty uses evaluation
stream or episode clusters. A weak read-out remains consistent with information loss, poor
conditioning, limited support, or insufficient probe capacity.


\paragraph{Execution and uncertainty.}
All three models use the same canonical padded frame batch and predictor batch. A strict parity
check against the public rollout function is bit-exact at the production shape, and model parameters
and BatchNorm statistics remain unchanged. This canonicalization was fixed after an initial smoke
parity failure and before any successful read-out fit; no tolerance, model weight, or metric was
changed. The analysis recomputes 1,080 metric rows and 720 paired contrasts from saved float32
predictions without refitting. Intervals resample the 100 evaluation environment streams jointly
across compared read-outs, with 2,000 replicates and seed 20260921. They condition on the fitted
model and probe. AUROC and average precision are reported as point estimates.

For the model-2, horizon-eight pre-reveal cohort, the forecast-versus-copy MAE difference is
$-0.63$ pixels with interval $[-1.41,0.13]$ under source-specific scaling. The transfer-versus-adapted
difference is $3.85$ pixels with interval $[3.48,4.22]$. Common scaling gives $-0.64$
$[-1.42,0.12]$ and $3.89$ $[3.52,4.25]$, respectively. These are descriptive conditional comparisons;
the strong adaptation result occurs in one of the three fitted models.


\subsection{Read-out errors and a post-hoc tracking audit}
The predeclared absolute-error table is retained below for every fitted model and both scaling
analyses. Its global training-mean baseline is evaluated on the selected pre-reveal slice; it is
not a slice-specific optimal constant. A separate audit of saved predictions examined correlations
and RMSE without refitting or modifying any prediction. The model-2 forecast-fitted probe has
Pearson correlation $0.246$ on the full wall-present cohort and $-0.175$ within the pre-reveal slice.
Its full-cohort RMSE is $18.65$ pixels versus $13.49$ for the global training-mean constant. Thus the
MAE improvement should not be described as reliable tracking or successful recovery of wall position.
A subsequent train-only constant-baseline audit uses the same horizon and cohort masks, with 1,509
training and 716 evaluation rows in the pre-reveal slice. The training-cohort mean position is
$29.433$ pixels and the median is $30$. Their evaluation MAEs are $4.769$ and $4.642$ pixels,
respectively, improving on every highlighted model-2 feature read-out in that slice. The full-cohort
training median has MAE $10.560$, compared with $9.112$ for the forecast-fitted probe, while the
latter's RMSE remains worse than the mean constant. This is an explicitly post-hoc interpretation
check, with no feature-probe refitting, clipping, latent rescaling, or alteration of the declared
metric outputs.

\begin{table}[ht]
\centering
\caption{\textbf{Read-out fitting changes measured forecast error.}
Eight-step wall-position MAE in pixels on 716 pre-reveal continuing contexts. R, F, and C denote
real future, forecast, and copied-current latents; arrows indicate fitting and evaluation domains.
Both train-only scaling analyses are shown. The global training-mean baseline has MAE $14.24$;
the post-hoc training-cohort mean and median have MAE $4.77$ and $4.64$.
The larger refitting gain occurs in model 2; all three models are retained.}
\label{tab:forecast-transfer}
\inputtable{forecast_transfer.tex}
\end{table}

\begin{figure}[t]
\centering
\paperfig{forecast_readout_v2.pdf}{1.0}
\caption{\textbf{Refitting changes forecast read-out error.} All three model seeds and both
predeclared normalization schemes are retained. The original baseline is the global training-label
mean evaluated on the selected slice; the added cohort-median constant is a post-hoc train-only
check. Lower slice MAE can reflect calibration as well as tracking;
the model-2 correlation audit and full-cohort RMSE prevent a reliable-position-recovery claim.}
\label{fig:forecast-transfer}
\end{figure}
